\section{Giới thiệu đề tài}
\label{sec:overview}
\begingroup
\setstretch{1.15}
\setlength{\parskip}{4pt}

\subsection{Bối cảnh và lý do chọn đề tài}
Sản phẩm thủ công, làng nghề, món ăn và địa điểm văn hóa là những điểm tiếp
cận cụ thể để tìm hiểu bản sắc địa phương. Một hình ảnh hay tên gọi giúp
nhận diện đối tượng, nhưng chưa thể hiện đầy đủ nguyên liệu, kỹ thuật,
tập quán sử dụng và câu chuyện của cộng đồng. Thông tin cần được đặt trong
bối cảnh địa phương và liên kết với nguồn tham khảo.

Người đọc có thể tiếp cận nội dung này qua các cổng du lịch, hồ sơ di sản
và bộ sưu tập trực tuyến. Vietnam Tourism tổ chức nội dung theo điểm đến
và trải nghiệm; Wonders of Vietnam của Google Arts \& Culture kết nối
câu chuyện văn hóa với tư liệu trực quan \cite{vietnam-culture,google-wonders}.
CultureMap VN tập trung vào trải nghiệm bắt đầu từ một tỉnh thành trên bản
đồ, tìm đến nội dung cụ thể, kiểm tra nguồn và hỏi thêm về bài đang đọc.

\subsection{Bài toán và hướng giải quyết}
\begin{itemize}
  \item \textbf{Khó bắt đầu khám phá:} người đọc có thể biết tên tỉnh nhưng
  chưa biết tên sản phẩm hoặc làng nghề. Hệ thống cung cấp bản đồ tỉnh thành
  song song với tìm kiếm theo tên địa điểm, tên bài và từ khóa.
  \item \textbf{Cần kiểm tra và sửa thông tin:} thành viên có thể gửi đề xuất
  bài mới hoặc báo thông tin sai lệch kèm minh chứng. Admin đối chiếu nguồn
  trước khi công bố bài hoặc cập nhật nội dung.
  \item \textbf{Muốn tìm hiểu thêm:} phần hỏi đáp sử dụng tư liệu của bài
  đang đọc để trả lời kèm căn cứ; thông báo khi chưa đủ thông tin.
\end{itemize}

Ví dụ, khách chọn một tỉnh có dữ liệu, mở bài về làng nghề và hỏi nguyên
liệu tạo ra sản phẩm. Nếu phát hiện tên gọi chưa chính xác, thành viên
chọn báo sai lệch và đính kèm minh chứng. Admin đối chiếu tài liệu và có
thể công bố một phiên bản đã sửa.

\subsection{Mục tiêu}
Mục tiêu tổng quát là xây dựng website khám phá văn hóa theo địa phương,
có dữ liệu biên soạn truy nguyên được nguồn và hỗ trợ hỏi đáp trong phạm vi
hẹp. Mục tiêu cụ thể gồm:
\begin{enumerate}
  \item Cho phép tìm kiếm, chọn tỉnh trên bản đồ và đọc nội dung công khai.
  \item Phân biệt khách vãng lai, thành viên và admin bằng quyền rõ ràng;
  cho phép thành viên gửi bài đề xuất, báo sai lệch và theo dõi kết quả.
  \item Hỗ trợ admin biên soạn, kiểm chứng nguồn, công bố bài và xử lý
  yêu cầu của thành viên. Lưu lịch sử phiên bản để đối chiếu thay đổi.
  \item Tích hợp RAG để hỏi đáp có căn cứ từ tư liệu bài đọc; thực hiện
  fine-tuning quy mô nhỏ và đánh giá trước/sau trên một nhiệm vụ phù hợp.
  \item Liên kết yêu cầu với use case, mô hình dữ liệu, thiết kế tương tác,
  giao diện và ca kiểm thử để kiểm chứng phần được hiện thực.
\end{enumerate}

\ReportNeedSpace{33mm}
\subsection{Đối tượng sử dụng}
\textbf{Khách vãng lai} khám phá, đọc nguồn và hỏi thêm mà không cần tài
khoản. \textbf{Thành viên} đã đăng nhập có thể gửi đề xuất bài mới, báo
sai lệch và theo dõi kết quả xử lý. \textbf{Admin} quản lý bài, nguồn,
phiên bản và kiểm chứng các yêu cầu. Nhóm thực hiện chuẩn bị dữ liệu
ban đầu để hệ thống có nội dung khi đưa vào sử dụng.

\subsection{Phạm vi đề tài}
\subsubsection{Phạm vi nội dung và địa lý}
Nội dung tập trung vào sản phẩm đặc trưng, nghề và làng nghề, ẩm thực, địa
điểm văn hóa. Mỗi bài có một tỉnh thành chính để điều hướng và một chủ đề
chính; những liên hệ với địa phương khác được giải thích trong bài, không
mang ý nghĩa khẳng định độc quyền văn hóa cho một tỉnh.

Bản đồ thể hiện Việt Nam và dự kiến có dữ liệu cho khoảng 3--5 tỉnh/thành
đại diện Bắc, Trung và Nam. Những tỉnh còn lại xuất hiện trên bản đồ nền
ở trạng thái bất hoạt. Tên tỉnh, mã liên kết bản đồ và mốc địa giới phải
thống nhất. Tên địa danh trong tài liệu cũ được giữ theo bối cảnh lịch sử.

\subsubsection{Các nhóm chức năng của hệ thống}
\begin{longtable}{|L{33mm}|L{111mm}|}
\caption{Các nhóm chức năng của phiên bản đồ án}\label{tab:scope}\\\hline
\rowcolor{gray!15}\textbf{Nhóm} & \textbf{Nội dung dự kiến}\\\hline
Khám phá văn hóa & Tìm kiếm, chọn tỉnh có dữ liệu trên bản đồ, lọc theo danh mục, đọc bài và nguồn.\\\hline
Nội dung và kiểm chứng & Thành viên gửi bài đề xuất, báo sai lệch và theo dõi kết quả. Admin đối chiếu nguồn, công bố hoặc cập nhật bài và lưu lịch sử phiên bản.\\\hline
AI hỗ trợ tìm hiểu & RAG theo tư liệu bài đọc; thử nghiệm fine-tuning quy mô nhỏ trên nhiệm vụ do nhóm thống nhất và đánh giá kết quả.\\\hline
\end{longtable}
Phạm vi hiện thực và yêu cầu chi tiết được trình bày ở Chương~\ref{sec:requirements}.

\subsubsection{Giới hạn}
Đặt tour, đặt vé, thanh toán, chỉ đường, đánh giá sao và mạng xã hội nằm
ngoài phạm vi. Quyền công bố và sửa bài chính thức thuộc admin. Địa phương
và danh mục được nạp sẵn. Lưu yêu thích và cá nhân hóa là chức năng tùy
chọn. Phần AI gồm RAG theo tư liệu đã kiểm chứng và fine-tuning quy mô nhỏ
trên mô hình có sẵn. Nhiệm vụ và cách đánh giá được trình bày ở
Chương~\ref{sec:ai}.

\subsection{Ràng buộc và tiêu chí chất lượng}
Nhóm gồm ba sinh viên; thiết kế ưu tiên ít thành phần vận hành và dữ liệu
có thể chuẩn bị thủ công. Nguồn văn hóa và lớp địa giới phải có xuất xứ,
điều kiện sử dụng và mốc dữ liệu. Yêu cầu, quyền, CSDL và hành vi hiện thực
phải nhất quán. Các mục tiêu chất lượng và cách kiểm chứng được trình bày ở
Chương~\ref{sec:requirements} và Chương~\ref{sec:evaluation}.

\subsection{Chuẩn bị nội dung và nguồn dữ liệu}
Nhóm biên soạn dữ liệu ban đầu từ cơ quan văn hóa, bảo tàng, hồ sơ di sản
hoặc tài liệu có đơn vị phát hành rõ. Bài giới thiệu ghi nguồn hỗ trợ và
xuất xứ hình ảnh; thành viên khác đọc đối chiếu dữ kiện trước khi admin
công khai. Đây là trách nhiệm trong nhóm, không tạo thêm vai trò phần mềm.
Các dữ kiện còn mâu thuẫn cần được làm rõ hoặc tạm chưa sử dụng. Nguồn lớp
bản đồ, điều kiện sử dụng và mốc địa giới được kiểm tra riêng để phù hợp
với danh mục địa phương của hệ thống.

\endgroup
